\section{调研结果}

\begin{frame}{调研范围与方法}
  \begin{columns}[T,onlytextwidth]
    \column{0.48\textwidth}
      \textbf{资料范围}
      \begin{enumerate}
        \item 学术论文与技术报告
        \item 官方文档与公开数据
        \item 代表性系统或案例
      \end{enumerate}
    \column{0.48\textwidth}
      \textbf{筛选标准}
      \begin{enumerate}
        \item 与研究问题直接相关
        \item 来源与方法可核查
        \item 结论具有可比性
      \end{enumerate}
  \end{columns}
\end{frame}

% 调研方法对比：演示双栏排版在真实内容中的用法
\begin{frame}{调研方法对比}
  \begin{columns}[T,onlytextwidth]
    \column{0.48\textwidth}
      \textbf{文献调研}
      \begin{itemize}
        \item 覆盖近三年的期刊与会议工作
        \item 以方法类别为主线归纳
        \item 便于比较指标口径
      \end{itemize}
    \column{0.48\textwidth}
      \textbf{系统调研}
      \begin{itemize}
        \item 选取有公开实现的代表性系统
        \item 记录实验环境与参数设置
        \item 便于复现与横向对比
      \end{itemize}
  \end{columns}
  \vspace{0.4em}
  \begin{block}{方法取舍}
    两类调研互补：文献给出趋势与共识，系统给出可核查的实现细节。
  \end{block}
\end{frame}

\begin{frame}{代表性工作对比}
  \small
  \begin{tabularx}{\linewidth}{@{}YYYY@{}}
    \toprule
    工作 & 研究对象 & 方法 & 主要结论 \\
    \midrule
    工作 A & 对象 A & 方法 A & 结论 A \\
    工作 B & 对象 B & 方法 B & 结论 B \\
    工作 C & 对象 C & 方法 C & 结论 C \\
    \bottomrule
  \end{tabularx}
  \sourceNote{来源：\parencite{example2026}（格式示例，请替换为真实文献）。}
\end{frame}

% 横向汇总表：适合把多个工作的指标放进同一张表对比。
% Y 列会自动换行并左对齐，行数多时优先精简列而不是缩小字号。
\begin{frame}{指标汇总}
  \small
  \begin{tabularx}{\linewidth}{@{}lYYYY@{}}
    \toprule
    工作 & 数据集 & 主要指标 & 结果 & 备注 \\
    \midrule
    工作 A & 数据集 A & 指标一 & 数值 A & 复现设置一致 \\
    工作 B & 数据集 B & 指标一 & 数值 B & 口径略有差异 \\
    工作 C & 数据集 C & 指标二 & 数值 C & 结果尚未复现 \\
    \bottomrule
  \end{tabularx}
  \sourceNote{来源：根据各工作原文整理；不同来源的评价口径可能不完全一致。}
\end{frame}
